
\input{../texlib/head}


\usepackage{tabularx}

\begin{document}
\ifdefined\BOOK
\else
\setcounter{chapter}{5}
\fi

\chapter*{Binary Trees.} 
\thispagestyle{firststyle}
\ctt{ Complete the Table, Continue the correctness of the minimal length. } 

Arrays, Linked Lists, Heaps, and Hash Tables are all ways to maintain collections. Each offers a different trade-off in query running times. For some of them, we imposed good worst-case bounds, while for others we can only compute expected costs (Hash Tables), detail are supplied in \Cref{tab:datastructs}. The obvious question is whether one could guarantee both fast access and update running times in the worst-case. 

In the following two lectures, we will see that using Balanced Binary Search Trees, one can provide both access and insertion/deletion in $O(\log n $. 

\begin{table}[h!]
\centering
\begin{tabularx}{\textwidth}{XXXXX}
\textbf{Data Structure} & \textbf{Access} & \textbf{Max/min} & \textbf{Insert/Delete} & \textbf{Max/min}  \\ \hline
Array                   & Data 2          & Data 3           & Data 3                 & $O(\log n)$          \\ 
Linked-List             & $O(n)$          & $O(n)$           & $O(1)$                 & $O(\log n)$          \\ 
Heap                    & $O(n)$          & $O(\log n)$      & $O(\log n) $           & $O(\log n)$          \\ 
Hash Table              & $O(n)$          & $O( )$           & $O( )$                 & $O(\log n)$          \\ 
\end{tabularx}
\caption{Running times respecting to data structure and task }
\label{tab:datastructs}
\end{table}


\section{Binary Trees and How to Encode Them}
We have already seen, in heaps, that organizing our data in a graph-like structure can offer a speed advantage. For future applications, and in particular for maintaining data in sorted order, we will have to encode our data using binary trees. In contrast to heaps, these trees may not be almost complete and also have to support pointer manipulations, specifically placing a binary tree as a left or right subtree of a given node. To enable this, we will have to treat the \textbf{right}, \textbf{left}, and \textbf{parent} as variables, in contrast to heaps where they are determined completely by the node index. We begin this section by stating definitions.



\begin{framed}
  \begin{definition}\hfill
    \begin{enumerate}
      \item Binary Tree: A tree in which any vertex has at most two children.
      \item A descendant of vertex $x$ is a vertex in the subtree whose root is $x$. A left descendant of vertex $x$ is either a vertex in the subtree whose root is the left child of $x$, or $x$'s left child.
      \item An ancestor of $x$ is a vertex to which $x$ belongs as a descendant.
      \item A leaf is a vertex without children.
      \item Height of vertex $x$ is the length of the longest simple path (without cycles) between $x$ and one of the leaves. The height of leaf is $0$.
      \item Height of the tree is the height of its root, which is usually denoted by $h$.
    \end{enumerate}
  \end{definition}
\end{framed}


    \begin{figure}[h]
      \centering
      \input{tikz-binary_1.tex}
      \caption{ Example of a binary tree with $12$ vertices. In blue, one of the longest paths from the root to one of the leaves, which is of length $4$. }
      \label{fig:<+label+>}
    \end{figure}


    \textbf{Question.} What are the maximal and minimal heights of binary trees over $n$ vertices?
    

    \textbf{Answer.}  The maximal height is $n-1$, and matches to trees in which any vertex has at most one child. The minimal height is $\log n$. Let's prove it by induction. 

    Denote by $H_{n}$ the minimal height of a binary tree with $n$ vertices. Observe that $H_{n+1} \ge H_{n}$, That's because adding a vertex to an $n$-size binary tree can't shorten any of the paths from root to leaves, particularly not the longest one.
    \begin{enumerate}
      \item Base. $n=1$, The tree consists of a single vertex, which is both the root and a leaf. Thus, the longest path is of length $0$, and also $\log 1 = 0$.
      \item Assumption. Assume the correctness of the claim for any $m < n$. 
      \item Step. Consider a binary tree over $n$ vertices with minimal hight. The longest path consists of an edge between the root and one of its children. Denote by $x$ the size of the root's left subtree, so the size of root's right subtree is $n-x$, Thus we get:
        \begin{equation*}
          \begin{split}
            H_{n} & \ge 1 + \max \left( H_{x}, H_{n-x} \right) 
          \end{split}
        \end{equation*}
By the monotonic property of $H_{n}$, we have that the above is minimized when both $H_{x}=H_{n-x}=H_{n/2}$ if $n$ is even, or $\max(x,n-x) = \lceil n/2 \rceil$ when $n$ is odd. So we get:
        \begin{equation*}
          \begin{split}
            \ge & 1 + H_{\lceil n/2 \rceil} \\ 
            = & 1 + \log(\lceil \frac{n}{2} \rceil ) = \log(n+1)
          \end{split}
        \end{equation*}
    \end{enumerate}


    \begin{figure}[h]
      \centering
      
\tikzset{%
  every neuron/.style={
    circle,
    draw,
    %fill=black,
    radius=0.1mm
  },
  neuron missing/.style={
    draw=none,
    scale=2,
    text height=0.25cm,
    execute at begin node=\color{black}$\vdots$
  },
}

\begin{center}

\begin{tikzpicture}%[ultra thick]




  \node [every  neuron ] (u-1) at ( 10 ,-0) { } ;
  \node [every  neuron ] (u-2) at ( 8 ,-1.5) { } ;
  \node [every  neuron ] (u-3) at ( 12 ,-1.5) { } ;
  \node [every  neuron ] (u-4) at ( 6.5 ,-3) { } ;
  \node [every  neuron ] (u-5) at ( 9.5 ,-3) { } ;
  \node [every  neuron ] (u-6) at ( 10.5 ,-3) { } ;
  \node [every  neuron ] (u-7) at ( 13.5 ,-3) { } ;
%  \node [every  neuron ] (u-8) at ( 0.5 ,-4.5) { } ;
%  \node [every  neuron ] (u-9) at ( 1.5 ,-4.5) { } ;
%  \node [every  neuron ] (u-10) at ( 2.5 ,-4.5) { } ;
%  \node [every  neuron ] (u-11) at ( 1.25 ,-6) { } ;
%  \node [every  neuron ] (u-12) at ( 1.75 ,-6) { } ;
%  \node [every  neuron ] (u-3) at ( 6 ,-3) { } ;
%\node (u-4) at ( 0.0 ,-6) { $ \vdots $  };
%
%\node [every  neuron ] (u-6) at ( 5 ,-6) { $l$  };
%\node [every  neuron ] (u-7) at ( 7 ,-6) { $r$  };
%\node  (u-8) at ( 5 ,-7) { $\vdots$  };
%\node  (u-9) at ( 7 ,-7) { $\vdots$  };
\draw [->] (u-1) -- (u-2);
\draw [->] (u-1) -- (u-3);
\draw [->] (u-2) -- (u-4);
\draw [->] (u-2) -- (u-5);
\draw [->] (u-3) -- (u-6);
\draw [->] (u-3) -- (u-7);
%\draw [->] (u-4) -- (u-8);
%\draw [->] (u-5) -- (u-9);
%\draw [->] (u-5) -- (u-10);
%\draw [->] (u-9) -- (u-11);
%\draw [->] (u-9) -- (u-12);





  \node [every  neuron ] (v-1) at ( 5 ,-0) { } ;
  \node [every  neuron ] (v-2) at ( 4 ,-1.5) { } ;
  \node [every  neuron ] (v-3) at ( 3 ,-3) { } ;
  \node [every  neuron ] (v-4) at ( 2 ,-4.5) { } ;
  \node [every  neuron ] (v-5) at ( 1 ,-6) { } ;
  \node [every  neuron ] (v-6) at ( 0 ,-7.5) { } ;
%  \node [every  neuron ] (v-7) at ( 6 ,-3) { } ;
%  \node [every  neuron ] (v-8) at ( 0.5 ,-4.5) { } ;
%  \node [every  neuron ] (v-9) at ( 1.5 ,-4.5) { } ;
%  \node [every  neuron ] (v-10) at ( 2.5 ,-4.5) { } ;
%  \node [every  neuron ] (v-11) at ( 1.25 ,-6) { } ;
%  \node [every  neuron ] (v-12) at ( 1.75 ,-6) { } ;
%  \node [every  neuron ] (v-3) at ( 6 ,-3) { } ;
%\node (v-4) at ( 0.0 ,-6) { $ \vdots $  };
%
%\node [every  neuron ] (v-6) at ( 5 ,-6) { $l$  };
%\node [every  neuron ] (v-7) at ( 7 ,-6) { $r$  };
%\node  (v-8) at ( 5 ,-7) { $\vdots$  };
%\node  (v-9) at ( 7 ,-7) { $\vdots$  };
\draw [->] (v-1) -- (v-2);
\draw [->] (v-2) -- (v-3);
\draw [->] (v-3) -- (v-4);
\draw [->] (v-4) -- (v-5);
\draw [->] (v-5) -- (v-6);
%\draw [->] (v-3) -- (v-7);
%\draw [->] (v-4) -- (v-8);
%\draw [->] (v-5) -- (v-9);
%\draw [->] (v-5) -- (v-10);
%\draw [->] (v-9) -- (v-11);
%\draw [->] (v-9) -- (v-12);
%%  node[right=6pt] {$y$} at
%\draw[decoration={brace,raise=6pt},decorate, thick]
%(8,-1.5) -- node[right=10pt, align=left] {  $n^{\prime}$-size binary tree, \\It's guaranteed that any node, \\beside the root, \\satisfies the heap inequality. } (8, -7.5);
%
%\draw[draw=gray, dashed] (4.5,-7.5) rectangle ++(3.5,6);
\end{tikzpicture} 


\end{center}

      \caption{ On the left, a binary tree with 6 vertices at height 5, while on the right, a binary tree with 7 vertices at height 2. }
      \label{fig:<+label+>}
    \end{figure}


    \subsection{Encoding}

We encode a binary tree by associating a field to each vertex $x$, representing its right, left children, and parent. We use the notation $x$.left to refer to the left child of $x$, although the physical implementation may differ conceptually. For example, the way binary trees are implemented in Cormen is through 4 arrays. The first stores the value of $x$, while the others store pointers of specific types. For instance, the array LEFT, where LEFT.$x$ stores the left child of $x$.


If nothing else has been mentioned, then we can assume that we can add additional fields to the vertices.


\section{Binary Search Trees.}
Now, we are ready to understand to store numbers (or keys), by binary trees. 

\begin{framed}
  \begin{definition} Binary search tree (BST) is a binary tree which any node $x$ of it: 
    \begin{enumerate}
      \item Contains a field key, storing a number $x$.key. 
      \item Any left descendant $y$ of $x$ satisfies $y$.key $\le$ $x$.key. 
      \item Any right descendant $y$ of $x$ satisfies $y$.key $\ge$ $x$.key. 
    \end{enumerate}
  \end{definition}
\end{framed}



    \begin{figure}[h]
      \centering
      \input{tikz-binary_1.tex}
      \caption{The heap after the exchange, at line (9). The induction assumption is applied over the right subtree, which is framed in gray.}
      \label{fig:<+label+>}
    \end{figure}

\textbf{Question.} Let $T$ be a binary search tree, Where are the minimum and maximum values of $T$? (most left and right nodes). 



\textbf{Answer.}  Consider the algorithm which gets down by taking only left turns until hitting edge, as in \Cref{alg:leftleft}, Guideline for correctness: denote by $z, r$ the minimal vertex by both value and height, or in other words, the highest vertex among the vertices with minimal value, and the root. Consider the path descending from $r$ to $z$. Suppose that while walking through the path, we take a right turn at vertex $y$. Then, by the BST property, $y \le z$. If $y < z$, then it contradicts $z$ being minimal by value. Otherwise, if $y = z$, then it contradicts $z$ being the highest vertex among the minimal vertices by value.

\begin{algorithm}
\SetAlgoLined
\KwData{$T$ - tree}
\KwResult{pointer to vertex with minimum value in $T$}
    \If{$T$ is empty}{
        \Return null
    }
    \If{$T$.left is empty}{
        \Return $T$
    }
    \Return Min($T$.left)
\caption{Min} 
\label{alg:leftleft}
\end{algorithm}
\textbf{Question.}  Given that the value $x$ is in $T$ and that it's unique, how do you find the vertex that stores it?



\textbf{Answer.} Consider algorithm \Cref{alg:search}.


\begin{algorithm}
  \label{alg:search}
\SetAlgoLined
\KwData{$T$ - tree, key - key to search for}
\KwResult{pointer to vertex with key equal to key}
    \If{$T$ is empty}{
        \Return null
    }
    \If{$T$.key equals key}{
        \Return $T$
    }
    \If{key is less than $T$.key}{
      \Return Search($T$.left, key)
    }
    \Else{
      \Return Search($T$.right, key)
    }

\caption{Search} 
\end{algorithm}

\begin{framed}
  \begin{definition}Let $T$ be a binary search tree, and let $x$ be a node belonging to it. 
    \begin{enumerate}
        \item The \textbf{predecessor} of $x$ will be defined as a vertex $y$ such that $y$.key $\leq x$.key and $y$.key is maximal among the nodes satisfying this condition. If we were to set the values of $T$ in sorted order, then the predecessor of $x$ would be located on its left.
          \item The \textbf{successor} of $x$ will be defined as $y$, where $x$ is the predecessor of $y$.
    \end{enumerate}
  \end{definition}
\end{framed}


\textbf{Question.}

\textbf{Answer.}


\begin{algorithm}
\SetAlgoLined
\KwData{$x$ - vertex}
\KwResult{pointer to predecessor of $x$}
    \If{$x$.left is not empty}{
      \Return Max($x$.left)\;
    }
    $y \leftarrow x$.parent\;
    \While{$y$ is not empty and $x$ is $y$.left}{
        $x \leftarrow y$\;
        $y \leftarrow y$.parent\;
    }
    \Return $y$\;
\caption{Predecessor}
\end{algorithm}



\input{../texlib/tail}


